

In this study we aimed to understand whether the natural regeneration of tropical secondary forests is influenced by the surrounding forest landscape comparing a dry and a wet tropical forest. Specifically, our objective was to analyse the impact of forest cover, forest connectivity, and forest successional stage in the surrounding forest landscape (1000~m radius) on the seed richness, seed dispersal mode and the seed guild of woody plants in the regenerating forests. 

%The following section elaborates on the results of this research in relation to our research questions: \\
%\textit{1) Do forest cover, forest successional stage and forest connectivity predict variation in woody perennial seed richness in naturally regenerating tropical forest during early secondary succession?}\\
%\textit{2) Do the relationships between forest factors and seed richness differ between wet and dry tropical forests?}.

\textbf{Seed richness}\\
The seed richness was hypothesized to increase with higher forest cover \citep{arroyo-rodriguezLandscapescaleForestCover2023a, PEREZCARDENAS2021118818, arroyo-rodriguezMultipleSuccessionalPathways2017}; higher connectivity \citep{huancanunezSeedrainsuccessionalFeedbacksWet2021, hordijkLandUseHistory2024, v.h.safarLandscapeOpennessHas2022a} and higher cover of late successional forest \citep{martinez-ramosNaturalForestRegeneration2016, goosemForestAgeIsolation2016} in the surrounding forest landscape in both dry and wet tropical forest. 

Forest type was found to be the dominant predictor of seed richness with wet forest having significantly higher seed richness than dry forest (fig. \ref{fig:stats} c)). This likely reflects the more favorable ecological conditions of wet tropical forest such as higher availability of water and longer growing season which results in larger species pool and faster increase in tree diversity \citep{souzadepaulaRoleSeedRain2023, hordijkLandUseHistory2024}.
Furthermore, seed richness is likely lower in dry forest due to common water stress in the environment which pushes for dependence on resprout reproduction \citep{souzadepaulaRoleSeedRain2023}. 

%\textit{(impact of limited nutrients in the soil after agri use?)}

While effects were only marginally significant ($p<0.1$), we found that a higher proportion of late successional forest in the surrounding forest landscape increased seed richness in the seed rain in both dry and wet tropical forests (fig. \ref{fig:stats} a)). This pattern suggests that older forests provide more diverse and abundant seed sources, consistent with studies showing that mature forests act as reservoirs of species diversity, also known as the "external ecological memory", and contribute substantially to the seed rain in regenerating stands \citep{martinez-ramosNaturalForestRegeneration2016, chazdonNaturalRegenerationTool2016, arroyo-rodriguezMultipleSuccessionalPathways2017}. Furthermore, these results also align with \cite{goosemForestAgeIsolation2016}, who found forest age to be the most important predictor of species richness. Though, it is important to mention that they analysed species richness of established plants as opposed to seed rain richness \citep{goosemForestAgeIsolation2016}. 
%Additionally, \cite{souzadepaulaRoleSeedRain2023} found that mature forests contain a richer soil seed bank, which serves as a valuable reservoir of propagules. These seeds remain dormant until favorable conditions trigger their germination, allowing them to grow into plants that can subsequently disperse seeds into regenerating forests. In this way, soil seed bank of late successional forest in the surrounding forest landscape can indirectly increase the number of seeds dispersed into the regenerating forest. 

In contrast, we did not find any significant relationship between seed richness and forest cover despite the positive relationship hypothesized by the "habitat amount hypothesis"  \citep{fahrig2013rethinking} and known literature \citep{arroyo-rodriguezLandscapescaleForestCover2023a, PEREZCARDENAS2021118818, arroyo-rodriguezMultipleSuccessionalPathways2017}. This might be due to the fact that our model controlled for forest type. Meaning that not for both wet and dry forest type, seed richness increases with forest cover. This theory is supported by the significant positive correlation found between forest cover and seed richness in the wet forest (fig. \ref{fig:cor}), and no such relationship found in the dry forest. A possible explanation is that in the dry forest, the higher proportion of abiotically dispersed seeds makes seed dispersal less dependent on forest cover, since the primary dispersal agent is wind. %Furthermore, the negative correlation between forest cover and seed richness in the dry forest suggests that larger forest cover surrounding regenerating sites in dry forest can even inhibit seed rain recovery. 
In contrast, wet forest is dominated by biotically dispersed seeds, whose dispersal depends more strongly on surrounding forest cover because it relies on animals as the main dispersal agents \citep{hordijkLandUseHistory2024}. 

Similarly to forest cover, we found no significant relationship between forest connectivity and seed richness. This could be due to several reasons. First, the spatial scale of our analysis may have influenced the results: we measured connectivity within a 1000~m radius around regenerating forests, which corresponds to the typical movement distance of animal seed dispersers. However, this might not fit the regional animal dispersal distance \citep{corlettSeedDispersalDistances2009}. Following the research of \cite{v.h.safarLandscapeOpennessHas2022a}, the scale at which animal dispersed seeds were impacted was more than 1600~m radius. In contrast, the same research showed that the scale at which small seeds, which are often abiotically dispersed, were impacted, was much smaller ($<400$~m) \citep{v.h.safarLandscapeOpennessHas2022a}. Second, many seeds may rely on abiotic dispersal agents such as wind, reducing their dependence on the interconnectedness of forest patches, which primarily affects biotically dispersed seeds \citep{dentUnitingNicheDifferentiation2021}. %This can be observed in the Spearman's correlation where abiotic dispersal has either very weak correlation with forest connectivity ($\rho = 0.09, p > 0.05$; dry forest; \ref{fig:cor}) or even negative correlation ($\rho = -0.31, p > 0.05$; wet forest; \ref{fig:cor}). 
This is also supported by the negative correlation present between seed richness and forest connectivity in the dry tropical forest, dominated by abiotic seed dispersal and no such correlation found in the wet tropical forest, dominated by biotic seed dispersal (fig. \ref{fig:cor}). 
Finally, forest connectivity is a complex variable, and there is an ongoing debate about the best way to measure it \citep{tischendorf2000should}. It is possible that the method we used did not fully capture the relevant aspects of patch interconnectedness. Furthermore, applying the mean of euclidean nearest neighbor distance metric when evaluating the forest connectivity might obscure the actual variation in connectivity. While the mean defines the average, it potentially does not capture the extreme distances which might be highly impactful for animal dispersers. An alternative metric to use, could be the degree of aggregation of forest patches \citep{he2000aggregation, mcgarigal2002fragstats}.  

%\textit{(Perhaps, there are more important factors impacting seed richness such as land use history \citep{hordijkLandUseHistory2024, jakovacRoleLanduseHistory2021}?)}


\textbf{Seed dispersal mode}\\
The biotic seed dispersal was hypothesized to increase with higher forest cover \citep{chazdonDriversBenefitsNatural2025, martinez-ramosNaturalForestRegeneration2016}, connectivity \citep{hordijkLandUseHistory2024, v.h.safarLandscapeOpennessHas2022a} and with larger cover of late successional forest \citep{chazdonDriversBenefitsNatural2025} in the surrounding forest landscape. In contrast, the abiotic seed dispersal was hypothesized to decrease with an increase of forest cover \citep{poorterComprehensiveFrameworkVegetation2024}, connectivity \citep{damschen2014fragmentation} and late successional forest cover \citep{poorterComprehensiveFrameworkVegetation2024}. These relationships were hypothesized to be true in both forest types. 

The distinction between dry and wet forest was significant with wet forest having higher proportion of biotically dispersed seeds and dry forest having higher proportion of abiotically dispersed seeds (fig. \ref{fig:stats} c)). These results 
align with the literature which states that the dry forest's woody flora tends to be dominated by abiotically dispersed species \citep{souzadepaulaRoleSeedRain2023}. Dry forests are characterized by seasonal water scarcity, which promotes decidious vegetation as a strategy to conserve water resources. This results in more open canopy facilitating wind dispersal. Furthermore, the lack of water resources limits the production of fleshy fruits during the dry period and hence results in a lower presence of feeding animal dipsersers \citep{hordijkLandUseHistory2024, lohbeckSuccessionalChangesFunctional2013}. Moreover, there are also human induced impacts promoting abiotic dispersal in the dry tropical forest such as defaunation \citep{souzadepaulaRoleSeedRain2023, chazdonDriversBenefitsNatural2025} of the environment and historical high intensity land use \citep{jakovacRoleLanduseHistory2021, hordijkLandUseHistory2024}. In contrast to the dry forest, wet tropical forest provide high water availability, which supports dense, evergreen canopy and year-round fleshy fruits. This dense canopy structure and continuous fruit supply maintains a constant presence and activity of animal dispersers \citep{poorterComprehensiveFrameworkVegetation2024, i.hordijkLandUseHistory2024}. 

We did not find any significant relationship between the surrounding forest landscape structural factors and the seed dispersal modes. There could be various reasons for the absence of a direct relationship between the various forest landscape factors and the seed dispersal modes. %One, and perhaps the most probable, is the scale (1000~m) at which this research was done. \cite{v.h.safarLandscapeOpennessHas2022a} found that biotically dispersed seeds might be affected by forest cover at even larger scale ($>1600~m$) whereas small seeds, which are often abiotically dispersed at smaller scale ($<400~m$) \citep{v.h.safarLandscapeOpennessHas2022a}. 
For example, \cite{huancanunezSeedrainsuccessionalFeedbacksWet2021} found no significant increase of biotically or abiotically dispersed seeds with forest age, however, another study did find an increase in vertebrate dispersed seeds and forest age \citep{tabarelli2002abiotic}, suggesting a potential knowledge gap. 
%\textit{Our non-significant trend might occur also due to the presence of bats, which are present in high number even in human modified landscape with intense land use history \citep{huancanunezSeedrainsuccessionalFeedbacksWet2021}.} 
Additionally, other factors such as land use history might play an important role in the proportion of biotically or abiotically dispersed seeds. For instance, intensive land use may have led to a potential extirpation of animal dispersers and depletion of soil seed banks, thereby reducing the overall seed rain  \citep{jakovacRoleLanduseHistory2021}. 



%\textit{(proportion does not take into account the total number, which would obviously differ between wet and dry forest due to longer season in the wet forest}

\textbf{Seed guild}\\
The pioneer and generalist guilds were hypothesized to decrease with forest cover \citep{lohbeckSuccessionalChangesFunctional2013, jakovacRoleLanduseHistory2021}, connectivity \citep{dentUnitingNicheDifferentiation2021, hordijkLandUseHistory2024} and late successional forest cover \citep{chazdonDriversBenefitsNatural2025}, whereas shade-tolerant guild was hypothesized to increase with forest cover \citep{hordijkLandUseHistory2024}, connectivity \citep{santo-silvaNatureSeedlingAssemblages2013} and late successional forest in both forest types \citep{huancanunezSeedrainsuccessionalFeedbacksWet2021}.  

There was a significant positive correlation detected between generalist guild and early successional forest in the dry forest (fig. \ref{fig:cor}), however, this relationship was not confirmed by the linear regression. The reason for this is, most probably, the fact that we controlled for forest type in the linear models. This can be explained by the slight positive correlation among these two variable in the wet forest, which was probably not a strong enough relationship to have a predictive character. However, the linear regression did show a significant decrease in generalist seed species as forest cover increased (fig. \ref{fig:stats}). Though, among all other variables there was no significant relationship detected.

Our result, considering a decrease in seeds belonging to generalist guild as the surrounding forest landscape moves toward a larger forest cover, is generally supported by literature \citep{santo-silvaNatureSeedlingAssemblages2013, poorterComprehensiveFrameworkVegetation2024}. Generalist species often benefit from open, sun exposed landscape which often arises from disturbed and fragmented landscapes. Therefore, with increasing forest cover, the edge effect formed by disturbance in the environment decreases, leaving more space for specialists and less space for the generalist species \citep{v.h.safarLandscapeOpennessHas2022a, santo-silvaNatureSeedlingAssemblages2013}.

The reason for the absent significant relationship between shade-tolerant species and the forest attributes might be the high frequency of plots with zero or very low proportions of shade-tolerant species in the regenerating plots. Additionally, some structural forest landscape factors might be more important for one tropical forest type than for another. This can be observed in the significant negative correlation between shade-tolerant species and early successional forest and no such relationship found in the wet tropical forest (fig. \ref{fig:cor}).  Similarly to seed richness, we theorise that the detection of significant relationship between pioneer species and forest landscape factors could have been hindered by the radius extent of the surrounding forest landscape.
Another reason for the lack of significance among these relationships might be a potential misclassification of the seed species. 

\subsection{Strengths and limitations}
This study holds numerous strengths such as the use of novel change detection algorithm \ac{AVOCADO} with an integration of remote sensing data facilitating a long-term landscape perspective, doing a comparative study between wet and dry tropical forest type and elaborate analysis of the surrounding forest landscape not only considering the forest cover, but also including other forest landscape attributes such as forest connectivity or fragmentation and the successional stage at which the surrounding forest is. However, this research is also accompanied by many limitations such as small spatial scale ($\sim$ 3~km\ x\ 4~km) causing overlapping  landscape surrounding the regenerating forest research plots, short sampling duration (1 year), spatial (30~m) resolution, land cover map accuracy (80.6\%) and seed trap area difference in wet vs dry forest ($0.79~\text{m}^2$, $0.49~\text{m}^2$).

Compared to other change detection models, \ac{AVOCADO} algorithm offers: 1) the potential to detect multiple cycles of disturbance and regrowth, providing up-to-date insights into forest landscape dynamics, 2) flexibility in defining how long a disturbance or regrowth must last, enabling location-specific detection, 3) adjustable sensitivity, letting users tailor the detection to the type of disturbance of interest and 4) reference phenology derived from the same ecosystem, which accounts for natural variability of the ecosystem. 

\ac{AVOCADO} was supplied with Landsat satellite imagery which enabled for a long-term (30 years) forest landscape analysis. Remote sensing data are a valuable addition to the research as historical field data is often scarce and interview based data are often unreliable \citep{decuyperContinuousMonitoringForest2022}. Additionally, this study includes a dry tropical forest in the analysis which is often under-represented in the studies on succession of tropical forests \citep{chazdonDriversBenefitsNatural2025}. Lastly, by evaluating multiple forest attributes — forest cover, forest connectivity, and forest successional stage — this study provides a multivariate perspective that better reflects the ecological complexity of the surrounding environment.

While \ac{AVOCADO} offers many advantages, it also entails certain limitations in context of our research. To ensure that only  deforestation disturbances (as opposed to drought disturbances) were detected, we set the \ac{RFD} threshold to 99\%. This threshold implies that pixels divergent more than 99\% from the reference phenology are classified as disturbance. Although such a strict \ac{RFD} threshold increases the risk of omitting some deforestation events, we did not consider this a major limitation. Given our focus on deforestation disturbances, it was preferred to miss a small number of deforested pixels rather than risk including drought-related disturbances in the analysis (see \href{https://github.com/Bianky/MSc-Thesis/blob/main/Analysis/avo_calibration.xlsx}{\ac{RFD} threshold scenarios} for details).

Additionally, to avoid misclassifying crops as regrowth, the regrowth threshold was set to 365 days, meaning that only regrowth persisting longer than one year was detected. Consequently, regrowth initiated in 2022 could not be identified, as no satellite data beyond 2022 were collected, preventing the detection of one-year-old forest. 

The spatial resolution of Landsat (30~m) may also be considered a limitation, particularly for identifying small-scale disturbances such as selective or forest-edge logging. However, since the aim of this study was to compare forest landscapes using a consistent satellite imagery source, this limitation was not regarded as significant. Additionally, Landsat provides a valuable long-term record of satellite imagery, offering a continuity that higher-resolution satellites cannot match \citep{kennedyBringingEcologicalView2014a}.

Another technical constraint was the accuracy of the land cover map (tree cover class user's accuracy = 80.6\%). Additionally, aggregating the original land cover map to 30~m resolution may have further reduced this accuracy. Nevertheless, as this represents one of the highest-quality maps currently available, further improvement was not feasible within the scope of this research \citep{zanagaESAWorldCover102022}.

The relatively small extent of our \ac{ROI} resulted in overlapping surrounding forest landscapes for the research plots. Furthermore, since the regenerating plots are located fairly close to each other (see fig. \ref{fig:studyarea}), using a 1000~m radius to characterize the surrounding forest landscapes led to substantial overlap. Such overlaps may introduce bias in the statistical analysis and limit the ability to draw conclusions regarding the effects of the localized forest landscape on individual regenerating plots.

Regarding the statistical analysis, cross-validation of the models was not performed. While cross-validation would enhance model robustness, implementing a research-specific approach that accounts for spatial similarity among surrounding forest landscapes was beyond the scope of this study.

Another limitation is the relatively short sampling period of one year, which may underestimate inter-annual variability. However, monthly seed collection across 40 research plots in two regions of Mexico is highly labor-intensive, limiting the feasibility of longer-term sampling.

Lastly, differences in the methodological design of the seed traps represent a potential limitation. In the dry tropical forest, the traps were rectangular with an area of $0.49~\text{m}^2$, whereas in the wet tropical forest, the traps were circular with an area of $0.79~\text{m}^2$. This difference could contribute to the different amounts of seeds collected in wet vs dry tropical forest. 


\subsection{Further research } 
For future research, we recommend investigating how different radii can affect the results when analyzing the forest landscape surrounding regenerating research plots. Using a reduced radius may minimize  landscape overlap and thereby increase the independence of research units \citep{v.h.safarLandscapeOpennessHas2022a}. However, larger radius might potentially have more significant impact on animal dispersal \citep{v.h.safarLandscapeOpennessHas2022a}. Moreover, smaller radii may better capture the dynamics of small, abiotically dispersed seeds, potentially yielding different insights into how the surrounding forest landscape influences forest regeneration \citep{jakovacRoleLanduseHistory2021}. Additionally, to account for the overlap of the landscapes in any radius scenario, we would suggest to include the degree of overlap in the statistical analysis following the approach of \cite{v.h.safarLandscapeOpennessHas2022a}. 

Additionally, we recommend to standardize for the size of the seed traps to increase the comparability between the dry and the wet forest. 

Regarding the forest connectivity metric, we recommend to experiment using different approaches of evaluating the forest connectivity as the current method might not have captured the landscape fragmentation effectively. An alternative could be using number of forest patches within given radius or the degree of aggregation which shows how closely to each other are forest patches located \citep{he2000aggregation, mcgarigal2002fragstats}.

In our current methodology, we control for the forest type, but do not perform separate analyses for each forest type. As \cite{lohbeckSuccessionalChangesFunctional2013} emphasize the importance of analysing different forest types, we recommend conducting analysis separately for wet and dry tropical forests to improve understanding of their distinct successional dynamics. 

Finally, to strengthen ecological perspective, we recommend extending the duration of the seed sampling period to capture inter-annual variation. Long-term monitoring is essential for identifying the ecological drivers of secondary succession and therefore, natural regeneration \citep{chazdonNaturalRegenerationTool2016}.

\subsection{Nature conservation recommendations}
Natural regeneration as a passive form of restoration can be a valuable approach to recover tropical forests. This study showed that natural regeneration in areas surrounded by late successional forest can enhance seed rain richness, which can thereby promote establishment of a higher diversity of woody plants and development of forest structure. Furthermore, wet tropical forest was characterized by a higher proportion of biotically dispersed seeds and might therefore require higher forest connectivity as it is essential for animal dispersal movement. In contrast, dry tropical forest was characterized by higher proportion of abiotically dispersed seeds and might therefore be more dependent on greater openness in the forest landscape as less densely vegetated areas facilitate abiotic dispersal. Hence, tailoring the conservation plan to the region characteristics can have a significant impact on the success of natural forest regeneration. With these findings we recommend that when aiming naturally regenerate degraded landscape it is important to have a high proportion of late successional forest in the surrounding forest landscape as it increases the seed rain richness. In the wet tropical forest, forest connectivity is of a high priority, as wet tropical ecosystem showed higher dominance by biotically dispersed seeds. In the dry tropical it is important to have open areas in the surrounding landscape as dry tropical ecosystem showed higher dominance by abiotically dispersed seeds.



%With the findings of this study we can conclude that considering the successional stage of the forest in the surrounding forest landscape and having region specific nature conservation plan is essential when aiming to naturally regenerate degraded tropical landscape.


%kikc of 
%transition - generalist decline
 